\section*{अध्याय \ref{ch:b3:supramolecular} --- अधिआण्विक रसायन}

\begin{solution}{exo:b3:supramolecular:1}
आयन--द्विध्रुव; $\pi$ संचयन; धनायन--$\pi$; \omterm{def:b3:supramolecular:interactions}{हैलोजन आबंध} (I से N); तीन हाइड्रोजन आबंध।
\end{solution}

\begin{solution}{exo:b3:supramolecular:2}
$\Delta G^\circ = -RT\ln K = -8.314 \times 298 \times \ln(\num{1.0e4}) = \qty{-22.8}{kJ/mol}$।
\end{solution}

\begin{solution}{exo:b3:supramolecular:3}
$K_{\ce{K+}}/K_{\ce{Na+}} = 10^{1.7} = 50$; अंतर पोटैशियम के पक्ष में $RT\ln 50 = \qty{9.7}{kJ/mol}$ है।
\end{solution}

\begin{solution}{exo:b3:supramolecular:4}
$n/(n + m) = 1/3$: दो \omterm{def:b3:supramolecular:supramolecular}{अतिथियों} के लिए एक \omterm{def:b3:supramolecular:supramolecular}{परपोषी}, $\mathrm{HG_2}$।
\end{solution}

\begin{solution}{exo:b3:supramolecular:5}
$H_0 + G_0 + 1/K = \qty{4.667e-3}{mol/L}$; $[\mathrm{HG}] = (4.667 - \sqrt{4.667^2 - 4 \times 2.0 \times 2.0})/2 \times 10^{-3} = \qty{1.13e-3}{mol/L}$: बद्ध
अंश 0.566, $\delta = 3.560 + 0.160 \times 0.566 = \qty{3.651}{ppm}$।
\end{solution}

\begin{solution}{exo:b3:supramolecular:6}
$\log K = 18.3 - 8.6 = 9.7$, $K = \num{5e9}$; $\Delta G^\circ = -RT\ln K = \qty{-55}{kJ/mol}$। छह Ni--N आबंध दोनों ओर
एक ही प्रकार के हैं; मुक्त अणुओं की संख्या चार से
सात हो जाती है: स्थानांतरीय एन्ट्रॉपी का लाभ अभिक्रिया को प्रेरित करता है।
\end{solution}

\begin{solution}{exo:b3:supramolecular:7}
$\Delta G^\circ = -RT\ln(\num{2.0e5}) = \qty{-30.2}{kJ/mol}$; $T\Delta S^\circ = \Delta H^\circ - \Delta G^\circ = -40 + 30.2 = \qty{-9.8}{kJ/mol}$;
$\Delta S^\circ = \qty{-33}{J.K^{-1}.mol^{-1}}$: एन्थैल्पी-प्रेरित बंधन, जिसकी आंशिक कीमत एन्ट्रॉपी में चुकाई जाती है।
\end{solution}

\begin{solution}{exo:b3:supramolecular:8}
कॉपर(I) दो द्विदंतुर फ़ेनैन्थ्रोलीन इकाइयों को बाँधता है और, चतुष्फलकीय होने के कारण, उन्हें
समकोण पर थामता है, एक को उस स्थान से पिरोकर जहाँ दूसरे का वलय
बंद होगा। दोनों शृंखलाओं को वलयों में बंद करने से वे अंतर्बद्ध हो जाते हैं। सायनाइड का
बड़ा आधिक्य, जो कॉपर(I) को कहीं अधिक प्रबलता से बाँधता है, धातु को हटाकर
मुक्त \omterm{def:b3:supramolecular:interlocked}{कैटेनेन} छोड़ देता है।
\end{solution}

\begin{solution}{exo:b3:supramolecular:9}
$k_c = \pi\Delta\nu/\sqrt2 = \pi \times 120/1.414 = \qty{267}{s^{-1}}$। आयरिंग:
$\Delta G^\ddagger = RT\ln(k_BT/hk_c) = 8.314 \times 300 \times \ln(\num{6.25e12}/267) = \qty{59.6}{kJ/mol}$।
\end{solution}

\begin{solution}{exo:b3:supramolecular:10}
ठोस औषध के आधिक्य में मुक्त औषध $S_0$ पर नियत रहती है। तब $[\mathrm{DC}] = KS_0[\mathrm C]$ और
$C_t = [\mathrm C] + [\mathrm{DC}] = [\mathrm C](1 + KS_0)$, अतः $[\mathrm{DC}] = KS_0C_t/(1 + KS_0)$ और कुल विलेयता $S_0 + [\mathrm{DC}]$ है। यहाँ
$KS_0 = 0.050$: $S = 0.10 + 0.050 \times 10/1.050 = \qty{0.58}{mmol/L}$, नैज मान का लगभग छह गुना।
\end{solution}

\begin{solution}{exo:b3:supramolecular:11}
जब $1/K \ll H_0, G_0$, तो $[\mathrm{HG}] \approx \bigl(H_0 + G_0 - |H_0 - G_0|\bigr)/2 = \min(H_0, G_0)$: बद्ध अंश एक तुल्यांक तक $G_0/H_0$ है,
फिर 1। वक्र अब $K$ पर निर्भर नहीं करता: पर्याप्त बड़ा कोई भी मान
रव के भीतर वही तीक्ष्ण विराम देता है, अतः आँकड़े केवल यह दिखाते हैं कि $K$
किसी न्यूनतम से अधिक है।
\end{solution}

\begin{solution}{exo:b3:supramolecular:12}
$k_{\mathrm{off}} = k_{\mathrm{on}}/K = \qty{1e9}{L.mol^{-1}.s^{-1}}/\qty{1500}{L.mol^{-1}} = \qty{6.7e5}{s^{-1}}$। विस्थापन अंतर
$\qty{0.160}{ppm}$, \qty{400}{MHz} पर \qty{64}{Hz} है; संगलन के लिए $k \approx \pi \times 64/\sqrt2 = \qty{140}{s^{-1}}$ चाहिए। विनिमय
हज़ारों गुना तेज़ है: एक औसत संकेत।
\end{solution}

\begin{solution}{pb:b3:supramolecular:1}
\textbf{1.} छह $\ce{-CH2CH2O-}$ इकाइयों का एक वलय; संकुल में $\ce{K+}$ केंद्र पर बैठता है, भीतर की ओर इंगित
छह ऑक्सीजनों द्वारा बँधा, $\ce{CH2}$ समूह बाहर।

\textbf{2.} $\ce{K+}$: $-RT\ln 10 \times 6.1 = \qty{-34.8}{kJ/mol}$; $\ce{Na+}$: \qty{-25.1}{kJ/mol}।

\textbf{3.} $10^{6.1 - 4.4} = 50$।

\textbf{4.} बड़ा $\ce{K+}$ (\qty{138}{pm}) वलय के सभी छह ऑक्सीजनों तक एक साथ पहुँचता है;
छोटा $\ce{Na+}$ (\qty{102}{pm}) नहीं पहुँच सकता, जब तक वलय मुड़ न जाए, जिसमें तनाव की कीमत लगती है।

\textbf{5.} इसका बाहरी भाग एक हाइड्रोकार्बन पृष्ठ है और इसका आवेश दबा हुआ है; $\ce{KCl}$ को
अपने आयनों का विलायकन चाहिए, जो टॉलूईन नहीं कर सकता।

\textbf{6.} जल $\ce{K+}$ का मेथेनॉल से कहीं बेहतर विलायकन करता है, और वलय में प्रवेश के लिए आयन को वह
जल छोड़ना पड़ता है।

\textbf{7.} $\ce{K+}(\text{aq}) + \ce{MnO4-}(\text{aq}) + \mathrm L(\text{कार्बनिक}) \rightleftharpoons \ce{[KL]+MnO4-}(\text{कार्बनिक})$।

\textbf{8.} $[\ce{K+}]_{\mathrm{aq}}$, \qty{0.10}{mol/L} पर बना रहता है: $y = \num{1.0e5} \times 0.10 \times (\num{1.0e-3} - y)^2 = \num{1.0e4}(\num{1.0e-3} - y)^2$।

\textbf{9.} $u = \num{1.0e-3} - y$ के साथ: $\num{1.0e4}u^2 + u - \num{1.0e-3} = 0$, $u = \qty{2.70e-4}{mol/L}$, $y = \qty{7.30e-4}{mol/L}$: 73\,\%
निष्कर्षित।

\textbf{10.} $y = \num{1.0e4}(\num{1.0e-3} - y)(\num{1.0e-2} - y)$ से $y = \qty{9.89e-4}{mol/L}$ मिलता है: 99\,\%।

\textbf{11.} $y = \num{1.0e4}(\num{1.0e-3} - y)(\num{1.0e-4} - y)$ से $y = \qty{9.0e-5}{mol/L}$ मिलता है: 9\,\%, क्राउन द्वारा सीमित,
जो 90\,\% भरा है।

\textbf{12.} परमैंगनेट आयन टॉलूईन में अपना बैंगनी रंग बनाए रखता है। वहाँ
वह अविलायकित, एक नग्न ऋणायन, है और जल की तुलना में कहीं तेज़ी से अभिक्रिया करता है, जहाँ एक
जलयोजन कोश उसे परिरक्षित करता है।

\textbf{13.} सोडियम आयन दोनों अनुनाद आवृत्तियों के अंतर से कहीं तेज़ी से
आता-जाता है: तीव्र विनिमय।

\textbf{14.} $\delta = \delta_{\mathrm{free}} + (\delta_{\mathrm{bound}} - \delta_{\mathrm{free}})[\mathrm{HG}]/H_0$, जहाँ $[\mathrm{HG}]$ यथार्थ समतापी से, $H_0 = \qty{2.0e-3}{mol/L}$,
$G_0 = (\text{तुल्यांक})\,H_0$, $\delta_{\mathrm{free}} = \qty{3.560}{ppm}$।

\textbf{15.} $K = (1.55 \pm 0.08) \times 10^3\ \unit{L.mol^{-1}}$, $\delta_{\mathrm{bound}} = \qty{3.719 \pm 0.001}{ppm}$।

\textbf{16.} वर्ग-माध्य-मूल अवशिष्ट \qty{0.0014}{ppm} है, पाठ्यांक
परिशुद्धता के बराबर, बिना किसी प्रवृत्ति के: 1:1 मॉडल उपयुक्त है।

\textbf{17.} बद्ध अंश $f$, $G_0 = fH_0 + f/(K(1 - f))$ पर: $f = 0.2$ के लिए 0.28 तुल्यांक और $f = 0.8$ के लिए 2.1।

\textbf{18.} $KH_0 = 10^{6.1} \times \num{2.0e-3} \approx 2500$: वक्र एक तुल्यांक पर तीक्ष्णता से विराम लेता
और केवल निम्न सीमा देता; निम्नतर सांद्रता पर कोई स्पर्धा प्रयोग या ऊष्मामिति
आवश्यक है।

\textbf{19.} टॉलूईन में, ऐल्कीन और परमैंगनेट के एक समांगी विलयन में।

\textbf{20.} अपचयित मैंगनीज़ कार्बनिक प्रावस्था छोड़ देता है (मैंगनीज़ डाइऑक्साइड के रूप में),
और अंतरापृष्ठ पर मुक्त हुआ क्राउन एक और $\ce{K+}$ और $\ce{MnO4-}$ उठा लेता है। प्रत्येक क्राउन
अणु परमैंगनेट को अनेक बार वहन करता है: उत्प्रेरकी मात्रा पर्याप्त है।

\textbf{21.} $K_{\mathrm{ex}}$, $[\ce{K+}]$ से गुणा होता है: ऊँची पोटैशियम सांद्रता निष्कर्षण को आगे धकेलती है।

\textbf{22.} चतुष्क अमोनियम लवण, जैसे टेट्राब्यूटिलअमोनियम या
बेंज़िलट्राइएथिलअमोनियम क्लोराइड, जिनका धनायन अपने-आप में वसारागी है।

\textbf{23.} बेंज़ीन ल्यूकेमिया उत्पन्न करता है; टॉलूईन यही काम कहीं कम
ख़तरे के साथ करता है।

\textbf{24.} \qty{1.0}{mmol/L} क्राउन के साथ 73\,\% परमैंगनेट
टॉलूईन में निष्कर्षित होता है।
\end{solution}
